\section{Learning and Retention}
\label{sec:learning-main}
\label{sec:monodromy-learning-main}

\subsection{Joint operation formation from terminal feedback}

For observed examples $(x,u_{1:L},y)$, the current learner minimizes terminal
negative log likelihood,
\begin{equation}
 \mathcal L(\theta)=-\frac1B\sum_{i=1}^{B}
    \log p_\theta(y_i\mid x_i,u_{i,1:L_i}).
\end{equation}
The initial representation, primitive matrices, caller, memory read/write
projections and readout all receive gradients through the executed sequence.
The learner receives neither correct calls nor intermediate states. Primitive
formation here means changes to executable learned state updates, rather than
assignment of examples to an already supplied answer table. It does not imply
that a primitive becomes a separately named or independently interpretable
algorithm. Such behavior requires execution and intervention evidence.

Soft execution evaluates every primitive and combines its residual according
to the learned call probabilities. This supplies terminal credit to the whole
bank. Backpropagation through time also differentiates the learned content
memory; state and memory reset for each example. The continuous development
runs use soft training. Their optimization budget, sampled minibatches and
validation protocol are fixed independently of test scores. Learning curves
and failed predecessors are retained rather than replacing low results with
an assumed successful configuration.

\subsection{Soft and hard-call execution}

Hard-call execution selects
$\widehat k_{t,s}=\arg\max_k\alpha_{t,s,k}$ and evaluates only
$\delta_{t,s,\widehat k_{t,s}}$. Encoding, state, memory, primitive functions
and the probabilistic readout remain continuous. We call this interface
\emph{native} to distinguish it from mixture execution; it is not a fully
discrete program or a one-hot classifier. Both interfaces use the same saved
weights, and both yield normalized terminal distributions.

An optional straight-through interface uses hard calls in the forward pass
and a dense soft-mixture surrogate in the backward pass. Writing
$\delta^{\mathrm{soft}}=\sum_k\alpha_k\delta_k$ and using
$\operatorname{sg}$ for stop-gradient, its residual is
\begin{equation}
 \delta^{\mathrm{ST}}=
 \operatorname{sg}(\delta_{\widehat k})+
 \delta^{\mathrm{soft}}-\operatorname{sg}(\delta^{\mathrm{soft}}).
\end{equation}
This is a biased surrogate for a discontinuous decision, not its exact
gradient. Soft and native trajectories can diverge because subsequent calls
and memory writes depend on their resulting states. Consequently, likelihood,
accuracy and inference cost must be measured separately; correctness of soft
execution does not certify a hard-call program.

\subsection{Retention as an empirical obligation}

Shared updates may change old encodings, primitive behavior, calls, memory
access and readout. The present continuous Block has no built-in theorem
protecting their composition. Its continual-learning policy must therefore
be evaluated against replay, freezing, modular adaptation and other applicable
controls with the same information. Extra examples, snapshots, auxiliary
state and maintenance are counted. Finite current-task success cannot replace
measurements during and after later learning.

Earlier implementations supply a distinct optional protection mechanism:
homotopic deformations of a based loop preserve its monodromy action, while
additional finite closure checks protect learned task bindings. Those
constructions and their counterexamples remain in
Appendices~\ref{sec:monodromy-principle} and~\ref{sec:prometheus-block}.
Their certificates apply to the specified finite execution models, not to the
unconstrained nonlinear primitives above. Monodromy is not required for the
current Primary candidate; its later inclusion would need an explicit
integration and a measured benefit beyond strong ordinary controls. Neither
this optional component nor an independent new learning rule is presumed
necessary for a useful architecture.
